\section{Results}
\label{sec:results}
We first describe and check our implementation
of momentum Wuppertal smearing. 
Then, in Sec.~\ref{sec:momtest} we optimize the smearing parameters
and test the effectiveness of the method. In Sec.~\ref{sec:boo}
we investigate whether introducing an additional Lorentz contraction
is advantageous, before determining pion and nucleon dispersion relations
up to momenta of $1.94\units{GeV}$ and $2.82\units{GeV}$, respectively,
in Sec.~\ref{sec:dispersion}.
\subsection{Implementation of the smearing}
We iterate the momentum Wuppertal smearing on our
$d+1=3+1$ dimensional gauge ensemble described in Sec.~\ref{sec:simulate},
employing spatially APE smeared~\cite{Falcioni:1984ei} gauge
transporters within Eq.~\eqref{fermsmear2}. These are iteratively
constructed as follows.
\begin{equation}
\label{ape}
U_{\vect{x},i}^{(n+1)}= P_{\textmd{SU}(3)}\!\!\left(U_{\vect{x},i}^{(n)}+\apec\sum_{|j|\neq i}
U_{\vect{x},j}^{(n)}U^{(n)}_{\vect{x}+\unitj,i}U^{(n)\dagger}_{\vect{x}+\uniti,j}\right),
\end{equation}
where $i\in\{1,2,3\}, j\in\{\pm 1,\pm 2,\pm 3\}$:
the sum is over the four spatial ``staples'' surrounding
$U_{\vect{x},i}$ where, again, we suppressed the time index as the
smearing is local in time. $P_{\textmd{SU}(3)}$ is a gauge covariant
projector onto the $\textmd{SU}(3)$
group, defined by maximizing $\re\tr[A^{\dagger} P_{\textmd{SU}(3)}(A)]$.
We iterate over the three
diagonal $\textmd{SU}(2)$ subgroups to achieve this. Other
projection possibilities can, e.g., be found in
Refs.~\cite{Morningstar:2003gk,Bali:2005fu}.
We iterate Eq.~\eqref{ape}
15 times, using the weight factor $\apec=2.5$.

Momentum Wuppertal smearing
Eq.~\eqref{fermsmear2} is implemented, multiplying the APE smeared links
for a given $\vect{k}$ value by phases,
$U_{\vect{x},j}\mapsto U_{\vect{x},j}e^{i\vect{k}\cdot\unitj}$,
and then iterating the usual Wuppertal smearing
Eq.~\eqref{fermsmear} on these links. Within this article we set $\wupc=0.25$
to obtain smooth smearing functions at tolerable iteration counts.
Starting from three $\delta$-functions (one for each source colour $b$) at
the spatial position $\vect{0}$, we can define ``wave functions'' for
three different source colours,
\begin{equation}
\psi_{(\vect{k})\vect{x}}^{ab}=\sum_{\vect{x}',a'}\left(\Phi_{(\vect{k})}^n\right)_{\vect{x}\vect{x}'}^{aa'}\delta_{\vect{x}'\vect{0}}^{a'b}\,,
\end{equation}
and the associated gauge invariant density:
\begin{equation}
\label{density}
\rho(\vect{x})=\frac{\sum_{ab}|\psi_{(\vect{k})\vect{x}}^{ab}|^2}
{a^3\sum_{\vect{x},ab}|\psi_{(\vect{k})\vect{x}}^{ab}|^2}\,.
\end{equation}
$\rho(\vect{x})$ does not carry any information relating to the
$\textmd{U}(1)$ phases. We therefore define
\begin{equation}
\label{phase}
\varphi(\vect{x})=\arg\left(\frac{\psi_{(\vect{k})\vect{x}}^{11}}{
\psi_{(\vect{0})\vect{x}}^{11}}\right)\in[0,2\pi)\,,
\end{equation}
as the phase of momentum smearing,
relative to the standard Wuppertal smearing. Above, we have singled
out one particular colour component but any diagonal component
will give the same phase function $\varphi(\vect{x})$.
In the case of a free configuration, i.e.\ employing trivial
links $U_{\vect{x},i}=\mathds{1}$, Eqs.~\eqref{density}--\eqref{phase}
simplify since $\psi_{(\vect{k})\vect{x}}^{11}=\psi_{(\vect{k})\vect{x}}^{22}
=\psi_{(\vect{k})\vect{x}}^{33}$ and colour off-diagonal elements vanish.
Moreover, $\varphi(\vect{x})=\arg(\psi_{(\vect{k})\vect{x}}^{11})$ in this case.

In Fig.~\ref{fig:smearing} we display $\rho(\vect{x})$ in
lattice units for the two dimensional cross section $x_3=0$ and
different settings. The colour encodes the phase $\varphi(\vect{x})$.
We
employ $n=200$ momentum Wuppertal smearing iterations with
$\wupc=0.25$ and set $\vect{K}=(1,1,0)=[L/(2\pi)]\vect{k}$,
i.e.\ $\vect{k}$ lies within the depicted plane and its modulus
is given as $|\vect{k}|\approx 0.77\units{GeV}$. Up to discretization and
finite volume effects, with this smearing we expect the free
case variance $\sigma^2\approx (6.3a)^2$ from Eq.~\eqref{eq:width}, i.e.\
$\rho(|\vect{x}|=6.3a)\approx e^{-1}\rho(0)$. 
In the top left panel we show the free case result from our iterative
smearing, which is consistent with this expectation. In the top right
panel of Fig.~\ref{fig:smearing} we repeat this on one time slice
of one of our gauge configurations, after having APE smeared
the gauge links. The resulting shape is slightly narrower but
otherwise indistinguishable from the free case and almost invariant
with respect to continuous
rotations. However, rotations take
place in colour space: plotting individual components of
$q^{ab}_{\vect{x}}$ (not shown) gives a less smooth behaviour. In
particular, the off-diagonal components do not vanish anymore.
The smoothness of the gauge invariant density $\rho(\vect{x})$
means that the differences relative to the free case
can be removed almost completely by a suitable gauge transformation.

In the
bottom left panel we apply momentum Wuppertal smearing to a time slice of the
original, not APE smeared gauge links. In this case the resulting
density is less symmetric and less smooth.  As one
may expect from mean field
arguments~\cite{Parisi:1980pe}, the average
smearing radius is somewhat reduced in a way consistent with multiplying
$\wupc$ by the fourth root of the average plaquette. Also in
this case, the $\textmd{U}(1)$ phase information
is intact. A comparison with
the top right panel of Fig.~\ref{fig:smearing}
clearly demonstrates the advantage of additional gauge link
smearing.

Finally, 
in the lower right panel we apply boosted momentum Wuppertal
smearing Eqs.~\eqref{eq:boostsmear}--\eqref{boostnorm},
using the APE smeared gauge links. We set
$\gamma=5.3$, which corresponds to the ratio
of the pion energy for a momentum
$\vect{p}\approx 2\vect{k}$ over the pion mass $m_{\pi}=295\units{MeV}$.
The smearing parameter is converted
according to Eq.~\eqref{boosttrans}. Indeed, the perpendicular shape
in the central region is basically unaltered relative to the top panels
while in the direction parallel to the boost the density
is contracted by the $\gamma$ factor. Due to the large
numerical value of this factor there are slight deviations from
the theoretical expectation but
these discretization related effects can be removed by reducing the
smearing parameter $\wupc'$ and increasing the iteration count $n$,
keeping $\sigma$ constant, see
Eqs.~\eqref{eq:width} and \eqref{boosttrans}.

\subsection{Optimization and test of momentum smearing}
\label{sec:momtest}
We now compute smeared-point and smeared-smeared
pion and nucleon two-point functions at different lattice momenta
Eq.~\eqref{quantize}, where we denote momentum and smearing
vectors in physical units as $\vect{p}$ and $\vect{k}$, respectively,
and integer component lattice momenta as
$\vect{P}=[L/(2\pi)]\vect{p}$. Note that $\vect{K}=[L/(2\pi)]\vect{k}$
does not need to have integer valued components. We (mostly)
restrict ourselves to
\begin{equation}
\label{eq:zeta}
\vect{K}=\zeta\vect{P}\,,
\end{equation}
where the naive expectation would be
$\zeta=1/2$ for the pion and $\zeta=1/3$ for the nucleon,
see Sec.~\ref{freefield}.
From the resulting two-point functions, we define effective
energies
\begin{equation}
\label{effen}
E_{H,\mathrm{eff}}(\vect{p},t+a/2)=
\frac{1}{a}\ln\left[\frac{C_H(\vect{p},t)}{C_H(\vect{p},t+a)}\right]\,,
\end{equation}
where $H\in\{\pi, N\}$. For the non-perturbatively improved
action that we use, which contains a clover term that
couples adjacent time slices, the meaningful range of 
$t$ values is $t\geq 2a$, i.e.\ we plot effective energies, starting
at $2.5a\approx 0.18\units{fm}$.

\begin{figure}
\includegraphics[width=0.48\textwidth]{zeta_values_mes_094.pdf}
\caption{Effective pion energies for $\vect{P}=(1,1,0)$, corresponding
to $|\vect{p}|\approx 0.94\units{GeV}$ and different ratios
$\zeta$, see Eq.~\protect\eqref{eq:zeta}.
The line is the expectation from
the continuum dispersion relation.
Symbols are shifted horizontally
to enhance the legibility.}
\label{fig:zeta}
\end{figure}

We realized different numbers of iteration counts $n$ both for momentum
and for conventional Wuppertal smearing. In both cases the best results
in terms of the ground state overlaps for pions and nucleons at different
momenta were obtained 
within the range $200\lesssim n\lesssim 400$. For the results we present here
we set $n=200$, $\varepsilon=0.25$, corresponding to
$\sigma\approx 6.3a\approx 0.45\units{fm}$, see Fig.~\ref{fig:smearing}.
We average the pion two-point function propagating in the forward
and backward time directions (folding). For the
nucleon we can only make use of the forward correlation function,
since we only employed the projector
$\mathbb{P}=\frac{1}{2}(\mathds{1}+\gamma_4)$,
see Eq.~\eqref{interpol}.

In Fig.~\ref{fig:zeta} we show effective energies from
smeared-smeared two-point functions for different values
of $\zeta$. In this case $\zeta=0.8$ (red squares)
appears to be the best choice, however the results are relatively
robust against increasing or decreasing this by 20\%.
In general,
at our pion mass $m_{\pi}\approx 295\units{MeV}$
and momenta up to $\sim 3\units{GeV}$, the optimal $\zeta$ values
came out to be $\zeta\approx 0.8>1/2$ for the pion and
$\zeta\approx 0.45>1/3$ for the nucleon. As discussed 
in Sec.~\ref{freefield}, we can only interpret
$\vect{k}$ as the momentum carried by a single
quark in the non-interacting case. Nevertheless, finding
values that are larger than the naive expectation,
rather than smaller, was somewhat unexpected.
Since the deviation from the free field case is not uniform
but bigger for the pion than for the nucleon, it would be
interesting to extend our study to non-Gaussian smearing functions.

\begin{figure}
  \includegraphics[width=0.48\textwidth]{plot_mes_094.pdf}
\caption{Effective pion energies Eq.~\protect\eqref{effen}
for the lattice momentum
$\vect{P}=(1,1,1)$, corresponding to $|\vect{p}|\approx 0.94\units{GeV}$.
In the case of momentum smearing (squares), we set
$\vect{K}=\zeta\vect{P}$ with $\zeta=0.8$.
Solid symbols correspond to smeared-smeared, open symbols
to smeared-point two-point functions. Some data points
are shifted horizontally to enhance the legibility.
The expectations from
the continuum and lattice dispersion relations can be found in
Eqs.~\protect\eqref{eq:disperse} and \protect\eqref{eq:dispersepi}, respectively.
Symbols are shifted horizontally
for better legibility.}
\label{fig:mes1}
\end{figure}

\begin{figure}
  \includegraphics[width=0.48\textwidth]{plot_mes_188.pdf}
\caption{The same as Fig.~\protect\ref{fig:mes1} for
$\vect{P}=(1,1,1)$, corresponding to $|\vect{p}|\approx 1.88\units{GeV}$.
The effective energies without momentum smearing cannot be determined,
due to prohibitively large errors and non-monotonous behaviour
of the central values of the respective two-point functions.}
\label{fig:mes2}
\end{figure}

\begin{figure}
  \includegraphics[width=0.48\textwidth]{plot_nuc_094.pdf}
\caption{The same as Fig.~\protect\ref{fig:mes1} for the nucleon
and $\zeta=0.5$. The horizontal lines correspond to the
expectations from the continuum and lattice dispersion relations
Eqs.~\protect\eqref{eq:disperse} and \protect\eqref{eq:disperseN}, respectively.}
\label{fig:nuc1}
\end{figure}

\begin{figure}
  \includegraphics[width=0.48\textwidth]{plot_nuc_188.pdf}
\caption{The same as Fig.~\protect\ref{fig:nuc1} for 
$\vect{P}=(2,2,2)$ ($\zeta=0.4$) and $\vect{P}=(3,3,3)$ ($\zeta=0.44$),
corresponding to
$|\vect{p}|\approx 1.88\units{GeV}$ and $|\vect{p}|\approx 2.82\units{GeV}$, respectively.
Again, the conventional Wuppertal smearing data are too noisy
to allow for the extraction of effective energies.}
\label{fig:nuc2}
\end{figure}

In Figs.~\ref{fig:mes1} and \ref{fig:mes2} we compare effective pion energies
from smeared-point (SP) and smeared-smeared (SS) correlation functions
with the expectations from the continuum and lattice dispersion relations
Eqs.~\eqref{eq:disperse} and \eqref{eq:dispersepi},
using the pion mass Eq.~\eqref{eq:masses} (horizontal lines). Note that
the SS effective energies should be monotonous functions of $t$ while
this need not be the case for SP energies. It is well known
that statistical errors of SP correlators are smaller than
in the SS case and we also confirm this. For the momentum smeared
two-point functions the data are consistent with plateaus
for $t\gtrsim 0.5\units{fm}$ and we find agreement with the expectations.
For $|\vect{p}|\approx 0.94\units{GeV}$ (Fig.~\ref{fig:mes1}) also the effective
energies from
conventionally smeared two-point functions are consistent with the
expectation, however, the errors are too large to allow for
quantitatively meaningful statements. For $|\vect{p}|\approx 1.88\units{GeV}$
(Fig.~\ref{fig:mes2}), within our statistics of one source position on
200 gauge configurations, it turned out to be impossible to obtain
effective energies without momentum smearing at all.

The same comparison is shown for the nucleon in
Figs.~\ref{fig:nuc1} and \ref{fig:nuc2}, where in Fig.~\ref{fig:nuc2}
we also include a momentum as high as $|\vect{p}|\approx 2.82\units{GeV}$.
In this case we show the lattice dispersion relation
Eq.~\eqref{eq:disperseN} of a free Wilson fermion with the
nucleon mass given in Eq.~\eqref{eq:masses}.
The statistical errors of the momentum smeared data again
are much smaller than for the non-momentum smeared cases.
Like for the pion, it was impossible within our statistics
to extract effective energies for momenta larger than $1.5\units{GeV}$.
At $|\vect{p}|\approx 0.94\units{GeV}$ the data agree with the expectation
for $t\gtrsim 0.6\units{fm}$ while for the higher momenta, where
the statistical errors are larger, the data are consistent
with plateaus starting at $t\gtrsim 0.45\units{fm}$.
For the high momenta shown in Fig.~\ref{fig:nuc2}
the data seem to prefer the continuum dispersion relation over the
lattice dispersion relation Eq.~\eqref{eq:disperseN}.

\begin{figure}
\includegraphics[width=0.48\textwidth]{extra_diff_dir_nuc.pdf}
\caption{The nucleon smeared-smeared effective energy with one
and the same momentum smearing $\vect{K}=(0.8,0.8,0.8)$,
optimized for $\vect{P}=(2,2,2)$, but for momenta pointing into
different directions.
Symbols are shifted horizontally
for better legibility.}
\label{fig:direction}
\end{figure}

It is clear from the results shown above that the gain of
using momentum smearing is tremendous. The only drawback, in addition
to the computational overhead from the smearing itself, is that
each value of the parameter $\vect{k}$ employed at the
source requires us to recompute the respective quark propagator.
A natural question therefore is whether it is possible
to efficiently realize several momenta $\vect{p}$ using one and
the same smearing vector $\vect{k}$.
We already know from Fig.~\ref{fig:zeta} that for
$\vect{p}\parallel \vect{k}$ the proportionality
constant $\zeta$ defined in Eq.~\eqref{eq:zeta} can be
varied by about 20\% without a significant deterioration
and by much more if one is willing to accept a compromise:
comparing Fig.~\ref{fig:zeta} with Fig.~\ref{fig:mes1}
reveals that even a much less than perfect momentum smearing
is a tremendous improvement over the conventional $\zeta=0$ case.
We may also ask whether it
is possible to (slightly) vary the direction of $\vect{p}$
relative to $\vect{k}$. This of course is potentially dangerous
since the interpolator used will not transform anymore
according to an irrep of the little group of $O_h$ (or its
double cover), associated to the momentum direction. However,
if for instance we are only interested in the ground state mass
of a spin-$0$ or spin-$1/2$ hadron this should not be a major problem.
Figure~\ref{fig:direction} demonstrates that to a certain extent varying
the momentum direction for a fixed $\vect{k}$
appears feasible too. In all the cases shown it is impossible to
extract any meaningful effective energies without momentum smearing.

\subsection{Test of boosted (momentum) smearing}
\label{sec:boo}
It has been suggested by two groups~\cite{Roberts:2012tp,DellaMorte:2012xc}
that introducing an anisotropy and thereby Lorentz boosting the smearing
function may improve the overlap of high momentum interpolators
with the respective hadronic ground states. We generalize
these ideas to off-axis momentum directions, also incorporating
phase factors, in the Appendix, see
Eqs.~\eqref{eq:boostsmear}--\eqref{boosttrans}.
Length contractions depend on the choice of coordinates
and in particular on time differences in the
moving frame relative to the rest frame. Since in Euclidean
spacetime all distances are spacelike and any real time
information is lost, it is not clear to us why spatial distances
should be subjected to a Lorentz boost.
Our numerical observations are negative.

\begin{figure}
\includegraphics[width=0.48\textwidth]{plot_SS_meson_boosted.pdf}
\caption{Effective smeared-smeared
pion energies for $\vect{P}=(3,0,0)$, corresponding
to $|\vect{p}|\approx 1.63\units{GeV}$:
momentum smearing with the optimized parameter $\zeta=0.8$
and boosted momentum
smearing Eq.~\protect\eqref{eq:boostsmear} for $\gamma\approx 5.6$
with $\zeta=0.5, 0.8, 0$.
Symbols are shifted horizontally
for better legibility.}
\label{fig:boost}
\end{figure}

In Fig.~\ref{fig:boost}, which is representative for our experiences,
we show effective pion energies for $\vect{P}=(3,0,0)$,
corresponding to $|\vect{p}|\approx 1.63\units{GeV}$ and a boost factor
$\gamma=E_{\pi}(p)/m_{\pi} \approx 5.6$, close to the one
that corresponds to the bottom right panel of Fig.~\ref{fig:smearing}
($\gamma=5.3$). In this case it was again not possible to extract
effective energies using the conventional Wuppertal smearing.
All Lorentz contracted (boosted) interpolators give results much
inferior to the one obtained using the unboosted momentum smearing.
This also holds for different contraction factors $1/\gamma$ (not shown).
At the same time the boost enhances the signal relative to
unboosted Wuppertal smearing, without the momentum phase factor.
In the figure we show boosted smearing results for different momentum shifts:
$\vect{k}=\vect{0}$, the naive expectation $\vect{k}=0.5\vect{p}$ and
the case that is close to optimal without the boost applied,
$\vect{k}=0.8\vect{p}$. We find
the effective energies and their uncertainties
are quite insensitive to the $\vect{k}$ value. 
This is not surprising since the
support of the smearing function in the direction of the momentum is
quite small. At the same time this small support
may explain why the boost outperforms
conventional Wuppertal smearing as broad wave functions are disfavoured
at high momenta, unless the $\vect{k}$ vector is introduced,
see Eq.~\eqref{eq:freec}.

In summary, substantially contracting the smearing function
in the direction of the momentum ameliorates the phase mismatch discussed
in this article. Therefore, some improvement over the conventional
isotropic smearing case can be achieved.
However, only momentum smearing correctly accounts for this effect
and we see no indication that injecting a momentum alters the
optimal shape of the modulus of the smearing function Eq.~\eqref{density}.

\subsection{Comparison with dispersion relations}
\label{sec:dispersion}
Our main aim here was to demonstrate the effectiveness of
momentum smearing. For this purpose it was sufficient to
consider only one source position on 200 individual
gauge configurations. The present state-of-the-art, however, is
to realize multiple sources on ten times as many
configurations.
In the near future we will compute a multitude
of physically interesting observables with enhanced statistics.
The masses shown in Eq.~\eqref{eq:masses} were already obtained
with high statistics and in Figs.~\ref{fig:zeta}--\ref{fig:boost}
we have compared effective energies against the
continuum and lattice dispersion relations
Eqs.~\protect\eqref{eq:disperse}--\protect\eqref{eq:disperseN},
using these values. 

\begin{figure}[ht]
\includegraphics[width=0.48\textwidth]{mes_disp_SS_8.pdf}
\caption{Pion energies for different lattice momenta.
in comparison to the continuum (solid curve) and
lattice (points connected by dotted lines) dispersion
relations Eqs.~\protect\eqref{eq:disperse} and
\protect\eqref{eq:dispersepi}.} 
\label{fig:disppi}
\end{figure}

\begin{figure}[ht]
\includegraphics[width=0.48\textwidth]{nuc_disp_SS_8.pdf}
\caption{Nucleon energies for different lattice momenta.
in comparison to the continuum (solid curve) and
lattice (points connected by dotted lines) dispersion
relations Eqs.~\protect\eqref{eq:disperse} and
\protect\eqref{eq:disperseN}.} 
\label{fig:dispN}
\end{figure}

In all cases the smeared-smeared effective energies from optimized
momentum smearing were in agreement with plateaus from
$t\geq t_{\min}=8.5a\approx 0.61\units{fm}$ onwards, where
$t=8.5a$ corresponds to the effective energy
obtained from the correlation function at
$8a$ and $9a$, see Eq.~\eqref{effen}. In many cases
$t_{\min}$ could be chosen smaller. For the moment being,
we conservatively approximate the energies by $E_H(\vect{p})\approx
E_{H,\mathrm{eff}}(\vect{p},t_{\min})$. The results as a function
of $\vect{p}$ are shown in Figs.~\ref{fig:disppi} and \ref{fig:dispN} and
compared to the dispersion relation expectations. We also display
results obtained with conventional smearing for small momenta where
this is possible. For the two-point functions studied here
the precision of the conventional results can be improved
at little computational overhead by averaging over (for the absolute momentum
values shown) six, eight or twelve equivalent directions.
We have not done this, to allow for a ``fair'' comparison of the
efficiency of the smearing methods. It is clear from the figures, however,
that the maximally possible error reduction, assuming
different momentum direction results to be
statistically uncorrelated, would not affect any of our conclusions.

We do not expect either parametrization shown in Figs.~\ref{fig:disppi} and
\ref{fig:dispN} to
perfectly describe the data as the lattice dispersion relations
are for point particles, assuming a particular form of the
effective Lagrangian. However, differences between the two functions
are indicative for the size of possible lattice effects.
While in the pion case differences between the parametrizations
are on the present level of statistics insignificant, the nucleon data
appear to be better described by the continuum dispersion relation.
In the near future we will further investigate this, increasing
our statistics and also employing a different smearing as described in
Sec.~\ref{sec:improve}.

